% TOP_PROBLEM: {"id": 27, "title": "Are there infinitely many perfect powers in the Fibonacci sequence?", "category_id": 1, "category": {"id": 1, "name": "number_theory", "display_name": "Number Theory", "description": "Properties of integers, prime numbers, Diophantine equations.", "slug": "number-theory", "order_index": 1, "created_at": "2026-07-31T15:26:25.670Z"}, "status": "open"}
% FIRST_POSTED: null
% ENTRY_KIND: statement_only
% Imported from: batch01.tex; UnsolvedMath ID 27
\documentclass[11pt]{article}
\usepackage[margin=25mm]{geometry}
\usepackage{fontspec}
\setmainfont{Latin Modern Roman}
\usepackage{amsmath,amssymb,mathtools}
\usepackage{unicode-math}
\setmathfont{Latin Modern Math}
\usepackage{xurl,hyperref}
\hypersetup{colorlinks=true,urlcolor=blue}
\setcounter{secnumdepth}{0}
\setlength{\parindent}{0pt}
\setlength{\parskip}{5pt}
\setlength{\emergencystretch}{3em}
\newcommand{\sref}[2]{\hyperlink{source:#1:#2}{[S#2]}}
\newcommand{\cataloguescope}[1]{\paragraph{Recorded target.}#1}
\begin{document}
% UnsolvedMath ID 27; source catalogue code NT-008
\setcounter{section}{26}
\section{Are there infinitely many perfect powers in the Fibonacci sequence?}\label{27}
{\small\sffamily UnsolvedMath ID 27 · Number Theory}
\subsection{Definitions and mathematical statement}

\paragraph{Shared notation.}$\mathbb N=\{1,2,\ldots\}$, and $\mathbb Z,\mathbb Q,\mathbb R,\mathbb C$ denote the integers, rationals, reals and complex numbers. Any other conventions are specified locally.


Use indices $n\in\{0,1,2,\ldots\}$ and define the Fibonacci numbers by
\[
 F_0=0,\qquad F_1=1,\qquad F_{n+2}=F_{n+1}+F_n.
\]
A nontrivial positive perfect power means an integer $a^b$ with integers $a\ge2$ and $b\ge2$. This is the convention $a,b>1$ in the catalogue's statement. Under this convention $1$ itself is excluded, even though it can be written $1^b$.
\paragraph{Statement recorded in the supplied source.}
Are there any nontrivial positive perfect powers in the Fibonacci sequence other than $8$ and $144$? Precisely, is it true that
\[
 \forall n\ge0\ \forall a,b\in\mathbb Z_{\ge2},\qquad
 F_n=a^b\ \Longrightarrow\ n\in\{6,12\}?
\]
Here $\mathbb Z_{\ge2}=\{2,3,\ldots\}$, and $F_6=8=2^3$, $F_{12}=144=12^2$. If bases $0$ and $1$ are also admitted, the corresponding classification is $F_n\in\{0,1,8,144\}$. \sref{27}{2}
The classification was proved by Bugeaud, Mignotte and Siksek; it also gives a negative answer to the title's question about infinitely many such values. \sref{27}{2}
\subsection{Short English statement}
Determine whether 8 and 144 are the only Fibonacci numbers that are powers with both base and exponent greater than one. The complete classification is already known; the exceptional values 0 and 1 appear only under the broader convention allowing those bases.
\subsection{Sources}
\begin{description}
\item[\hypertarget{source:27:1}{S1}] ULAM.AI and the UnsolvedMath Contributors, UnsolvedMath: A Curated Collection of Open Mathematics Problems, record 27 (NT-008). CC BY 4.0. \url{https://huggingface.co/datasets/ulamai/UnsolvedMath}.
\item[\hypertarget{source:27:2}{S2}] Yann Bugeaud, Maurice Mignotte and Samir Siksek, Classical and Modular Approaches to Exponential Diophantine Equations I. Fibonacci and Lucas Perfect Powers, Annals of Mathematics 163 (2006), 969--1018; arXiv:math/0403046, Theorem 1, page 2. \url{https://arxiv.org/pdf/math/0403046}.
\end{description}
\label{end:27}
\end{document}
